\renewcommand*{\authorOfNextLines}{Helmut Quinz}
\begin{quote}
	\gqqEmph{Wenn es offensichtlich ist, dass die Ziele nicht erreicht werden können, sollten die Handlungsschritte angepasst werden, nicht die Ziele.}\\
	\small{Konfuzius}
\end{quote}
Die Umsetzung des Projektes ist von der eigentlichen Projektarbeit geprägt. 
Sie werden mit hoher Wahrscheinlichkeit feststellen das manche Tätigkeiten sich nicht wie geplant verlaufen.
Einige Aufgaben lassen sich schneller lösen als vorgesehen, andere wehren sich gegen ihre Erfüllung.
Werden Sie nicht nachlässig sollte sich ersteres zu Beginn des Projektes häufen. 
Erfahrungsgemäß hält sich beides über das Projekt gesehen meist die Waage.\\
Stellen Sie sicher, dass alle am Projekt beteiligten (Team-Mitglieder, Betreuer, auftraggebende Unternehmen) zu jedem Zeitpunkt über den aktuellen Projektfortschritt informiert sind.
Das zwingt Sie automatisch zu einer Überwachung des Projekts und gibt den \gqq{Kunden} die Möglichkeit rechtzeitig den Projektfortschritt kritisch zu hinterfragen.\\

Dokumentieren Sie die Ereignisse, Erkenntnisse, Fortschritte und getroffenen Entscheidungen gewissenhaft. 
Dies erleichtert nicht nur die Nachverfolgung, sondern unterstützt auch die spätere Verfassung Ihrer offiziellen Projektdokumentation.

\paragraph{Wichtig:} Sie werden im Herbst von der Abteilungsleitung aufgefordert, Ihr Projekt und dessen Fortschritt zu präsentieren. 
Dies dient vor allem der Projektüberwachung.
Sie müssen jedoch davon ausgehen, dass diese Präsentation bei der Beurteilung der Arbeit berücksichtigt wird.